\section{Verification, limitations, and interpretation}
\label{sec:verification}
\subsection{What is verified}
Separate AI-assisted implementations reconstruct the equations without
importing the generating model. The exact logarithmic test reproduces
all 14,336 archived corner classifications; 11,002 tests check both sides
of finite positive productivity thresholds. The stored optimized shares
also pass objective, feasibility and concavity-gap checks. Scalar and
joint optimizations check the curvature examples, production alternative,
pure-computation budget solution and two-period extension; direct
accounting checks the workload examples, including the memory failure
beyond the stated horizon. Complete proofs establish the mathematical
claims; these numerical tests check implementations and examples, not
extraterrestrial premises. Supplementary Material sections~S1--S3 provide
certificate details, executable commands, numerical outputs and the
seed-versus-\texttt{rng} comparison. Frozen inputs, environment versions
and integrity manifests accompany the source.

\subsection{Limitations that change the conclusion}
\label{sec:limitations}
\paragraph{Physical support and recursive interpretation.}
The workload frontier is a homogeneous service model, not an engineered
architecture or a test of consciousness. Efficiency or fidelity may
plateau, stored state may dominate memory, and digital population growth
may restore scarcity. Per-agent resources increase only if useful
computational growth outpaces agent growth. The fixed-budget joint witness
chooses rather than derives its productivity path. No engineering model
here establishes indefinite growth of every effective support capacity.
Recursion is neither necessary nor sufficient for inward allocation.

\paragraph{Preferences, production and mathematical scope.}
The objective represents a specified decision process, not a shared
utility function of every inhabitant. Divisible investment, positive
prices, differentiability and concavity are substantive restrictions.
The local logarithm and external production functions are not empirically
identified; \cref{lem:production} separates corner robustness from
interior-allocation sensitivity. The value-curvature parameter also
changes the conclusion, as \cref{prop:curvature} shows. Small residual
utility weights do not certify physical safety or remove the value of
geographically independent backups. Fixed project sizes, start-up
thresholds and increasing returns can invalidate a local test. For example,
\begin{equation}
 U(x)=\ln\{1+2\ln(2-x)+0.1x\}+3x^2,\qquad 0\leq x\leq1,
 \label{eq:nonconcave}
\end{equation}
with $z=1-x$, has $U'(0)\simeq-0.3772$, yet
$U(1)\simeq3.0953>U(0)\simeq0.8697$. This violates the concavity
assumption rather than \cref{prop:corner}. Under added hard constraints,
necessity of the corner inequalities additionally requires the proposed
small budget-transfer directions to remain feasible.

\paragraph{Persistence and future opportunities.}
Proposition~\ref{prop:persistence} characterizes a sequence of static
optima, not an infinite-horizon policy. Its conditions are neither
necessary nor sufficient, in general, for an all-inward optimal policy
in a dynamic model: section~\ref{sec:dynamic} gives both reversals in two periods and
explains their extension when the later continuation is independent.
A substantive infinite-horizon criterion would require specified state
dynamics, feasible policies, and a well-defined discounted or long-run
objective. With logarithmic value, fixed positive residual slopes and bounded prices
exclude an all-inward path when budgets become unbounded. Residual decay
or sufficiently rising relative prices must satisfy the exact productivity
condition as well; simply observing that $S_j(t)\to0$ is insufficient.
The two-period model includes both local and external future resource
production, but not a full ecology of accumulated capital, strategic
competition, irreversible settlement, survival or real options.

\paragraph{Access and physical assumptions.}
The access-cost map is a sensitivity family, not a propulsion model.
Proposed warp geometries do not establish operational faster-than-light
travel \citep{Alcubierre1994}; energy-condition and causality objections
remain relevant \citep{Everett1996,Santiago2022}. The counterfactual
limiting cases assume consistent scheduling without usable travel into
one's own past. A rapid material courier may also carry information.
Faster access can increase both external opportunities and arrival
exposure, not merely reduce travel time.

\paragraph{Inference about populations.}
The synthetic design is not an astronomical prior. Neither its counts nor
variation across scrambles estimates the prevalence of inward civilizations.
Allocation does not identify power use, spectra, duty cycles, autonomous
expansionist exposure or arrival probabilities. Survey nulls are therefore
not selective confirmation of \CITH. Empirical application would require
independent evidence for relevant resource and preference parameters,
a propagation model and comparative observational likelihoods.

